\chapter{Leveraging Defects Lifecycle for Labeling Defective Classes}

\section{Contesto e obiettivo del lavoro}

\subsection{Defect prediction}

La \textbf{defect prediction} ha l'obiettivo di identificare gli artefatti software, ad esempio classi o file, che hanno un'alta probabilità di manifestare un difetto. Lo scopo principale è ridurre i costi di testing e di code review, permettendo agli sviluppatori di concentrare gli sforzi sulle parti più critiche del codice.

\subsection{Importanza del dataset}

Affinché un modello di predizione sia affidabile, è fondamentale costruire \textbf{dataset di difetti di alta qualità}. L'intero lavoro si concentra infatti non tanto sul classificatore in sé, quanto sulla fase precedente, cioè la \textbf{costruzione del dataset} e in particolare il \textbf{labeling} delle classi come difettose o non difettose.

Se il labeling è inaccurato, il modello apprenderà relazioni sbagliate tra metriche e defectiveness. Per questo motivo, la qualità del dataset non è un dettaglio secondario, ma una condizione necessaria per ottenere risultati credibili.

\begin{notaprof}
Ci troviamo nella fase iniziale di \textbf{misurazione} e creazione del dataset. Prima ancora di addestrare un modello di Machine Learning, bisogna capire come etichettare correttamente le classi che in passato sono state davvero affette da difetti.
\end{notaprof}

\subsection{Labeling delle classi difettose}

Le slide presentano due strategie principali per etichettare le classi come difettose o non difettose:

\begin{itemize}
    \item usare \textbf{SZZ};
    \item usare la \textbf{Earliest Affected Version (AV) in JIRA}, quando disponibile.
\end{itemize}

Il problema è che entrambe queste fonti possono essere incomplete o inaccurate. Da qui nasce la necessità di studiare metodi automatici alternativi per ricostruire il ciclo di vita del difetto.

\subsection{Obiettivo generale dello studio}

Lo studio ha un duplice obiettivo.

Da un lato, vuole capire \textbf{quanto siano disponibili e affidabili le informazioni di AV} inserite dagli sviluppatori nei ticket Jira. Dall'altro, propone un nuovo metodo automatizzato, chiamato \textbf{Proportion}, basato sull'idea che i difetti abbiano un \textbf{life-cycle relativamente stabile}.

In sintesi, il lavoro vuole:

\begin{enumerate}
    \item misurare quanto sia realistico affidarsi ai dati di AV presenti in Jira;
    \item proporre un metodo alternativo quando l'AV manca o è incoerente;
    \item confrontare l'accuratezza del nuovo metodo con quella di diverse varianti di \textbf{SZZ}.
\end{enumerate}

\section{Concetti fondamentali}

\subsection{SZZ}

\textbf{SZZ} è l'algoritmo classico per identificare il momento in cui un difetto è stato introdotto nel codice. A partire da un \textbf{fix commit}, SZZ usa il meccanismo di annotation del versioning system, per esempio \texttt{git blame}, per risalire all'ultima modifica precedente delle linee corrette.

\subsection{Earliest AV in JIRA}

L'approccio \textbf{Earliest AV in JIRA} consiste nell'usare la prima \textbf{Affected Version} riportata manualmente dagli sviluppatori nel sistema di tracciamento, ad esempio Jira. Se questa informazione fosse sempre presente e consistente, costituirebbe un riferimento molto utile per il labeling.

\subsection{Defect lifecycle}

L'idea centrale del paper è che ogni difetto attraversi un \textbf{ciclo di vita} che può essere descritto in termini di versioni del software. Questo ciclo di vita collega il momento in cui il difetto viene introdotto, quello in cui viene osservato e quello in cui viene definitivamente corretto.

\subsection{Significato di IV, OV, FV e AV}

Per descrivere il life-cycle del difetto, il paper usa quattro nozioni fondamentali.

\begin{itemize}
    \item \textbf{IV (Injected Version)}: è la versione in cui il difetto viene introdotto nel codice.
    \begin{notaprof}
    È la prima versione in cui il bug esiste davvero nel sistema, anche se magari nessuno se n'è ancora accorto.
    \end{notaprof}

    \item \textbf{OV (Opening Version)}: è la versione in cui viene aperto il ticket del bug.
    \begin{notaprof}
    È il momento in cui la \emph{failure}, cioè il malfunzionamento osservabile, viene riconosciuta e registrata.
    \end{notaprof}

    \item \textbf{FV (Fixed Version)}: è la prima versione rilasciata dopo il \emph{fix commit}.
    \begin{notaprof}
    È la prima release in cui quel bug non dovrebbe più comparire.
    \end{notaprof}

    \item \textbf{AV (Affected Versions)}: sono tutte le versioni affette dal difetto, cioè l'intervallo di versioni comprese tra il momento di introduzione e quello di correzione.
\end{itemize}

In termini concettuali, le versioni difettose sono quelle comprese nell'intervallo \([IV, FV)\). Questo significa che la \textbf{Fixed Version} non è più affetta dal bug, mentre tutte le versioni precedenti a partire dalla \textbf{Injected Version} lo sono.

\section{SZZ: funzionamento e limiti}

\subsection{Come funziona SZZ}

SZZ individua i commit che hanno introdotto il bug tracciando la storia del codice sorgente. L'algoritmo parte da un \textbf{fix commit}, cioè da un commit che sappiamo correggere un bug, e usa quella correzione come indizio per cercare all'indietro il possibile \textbf{bug-introducing commit}.

\subsection{Uso di \texttt{git blame} / annotation mechanism}

Operativamente, SZZ prende le linee di codice modificate da un \textbf{fix commit} e usa \texttt{git blame} per risalire nel tempo, identificando il commit precedente che aveva introdotto o modificato quelle stesse righe.

L'assunzione di fondo è che il bug sia stato introdotto proprio nelle linee poi corrette dal fix. Quando questa assunzione è vera, SZZ può fornire una buona stima del momento di introduzione del difetto.

\subsection{Bug introducing commit}

Il commit individuato tramite backward search viene considerato il \textbf{bug-introducing commit}. Da quel momento in poi, la classe o il file coinvolto vengono etichettati come difettosi.

\begin{notaprof}
In pratica, il commit che ha scritto la forma sbagliata del codice viene trattato come l'origine storica del bug. È un'inferenza plausibile, ma non sempre perfetta.
\end{notaprof}

\subsection{Limiti del metodo}

Le slide chiariscono che SZZ ha limiti strutturali importanti.

\begin{itemize}
    \item \textbf{Non è perfettamente accurato}: il commit trovato da SZZ è una stima, non una verità certa.
    \item \textbf{Non riesce a identificare bene i regression bugs}: un codice può diventare errato solo in seguito a cambiamenti successivi del sistema, anche se in passato era corretto.
    \item \textbf{Non funziona bene quando il fix consiste solo in aggiunta di codice}: se il bug viene corretto introducendo codice nuovo, senza modificare righe già esistenti, SZZ non ha una traccia storica affidabile su cui applicare \texttt{git blame}.
\end{itemize}

\begin{notaprof}
SZZ può generare sia \textbf{falsi positivi} sia \textbf{falsi negativi}. Produce falsi positivi quando va troppo indietro nel tempo e attribuisce il difetto a linee che storicamente erano corrette; produce falsi negativi quando non riesce a risalire al vero punto di introduzione del bug.
\end{notaprof}

\section{Aim dello studio e Research Questions}

\subsection{Idea di stable life-cycle dei defect}

L'intuizione centrale dello studio è che i difetti abbiano un \textbf{ciclo di vita stabile}. In altre parole, pur variando da bug a bug, sembrerebbe esistere una relazione abbastanza regolare tra:

\begin{itemize}
    \item il numero di versioni che intercorrono tra introduzione del difetto e sua scoperta;
    \item il numero di versioni che intercorrono tra scoperta del difetto e sua correzione.
\end{itemize}

Questa ipotesi motiva il metodo \textbf{Proportion}: se tale proporzione è abbastanza stabile, allora è possibile stimare automaticamente la posizione dell'IV anche quando l'AV esplicita manca.

\subsection{Research Questions}

Le domande di ricerca dello studio sono quattro:

\begin{itemize}
    \item \textbf{RQ1}: le \textbf{Affected Versions (AV)} sono disponibili e consistenti in Jira?
    \item \textbf{RQ2}: i vari metodi, ad esempio SZZ e Proportion, hanno accuratezze diverse nell'etichettare le \textbf{Affected Versions}?
    \item \textbf{RQ3}: i metodi hanno accuratezze diverse nell'etichettare le \textbf{singole classi difettose}?
    \item \textbf{RQ4}: l'uso di metodi diversi porta all'identificazione di \textbf{feature importanti} diverse nei modelli di predizione?
\end{itemize}

In questo capitolo ci fermiamo a \textbf{RQ1} e \textbf{RQ2}.

\section{RQ1: Is the AV available and consistent?}

\subsection{Rationale}

La prima domanda vuole verificare se il campo \textbf{Affected Version} sia effettivamente utilizzabile come fonte affidabile per il labeling. Ci sono infatti due problemi distinti:

\begin{itemize}
    \item \textbf{availability}: gli sviluppatori potrebbero non compilare affatto il campo AV;
    \item \textbf{consistency}: il campo potrebbe essere compilato, ma in modo incoerente rispetto all'ordine temporale delle versioni.
\end{itemize}

\begin{notaprof}
Per esempio, se la versione più vecchia indicata come affetta risulta successiva all'apertura del ticket, allora quella AV non è coerente e non può essere usata come ground truth affidabile.
\end{notaprof}

\subsection{Design e selezione di progetti e bug}

Per RQ1 sono stati analizzati \textbf{212 progetti Apache} collegati sia a \textbf{Git} sia a \textbf{Jira}. Sono stati considerati solo gli issue report che soddisfacevano simultaneamente i seguenti criteri:

\begin{itemize}
    \item \texttt{Type == Bug};
    \item \texttt{Status == Closed} oppure \texttt{Resolved};
    \item \texttt{Resolution == Fixed}.
\end{itemize}

Inoltre, sono stati esclusi:

\begin{itemize}
    \item i difetti che non avevano un \textbf{Git-related fix commit};
    \item i difetti che \textbf{non erano post-release}.
\end{itemize}

Dal punto di vista sperimentale:

\begin{itemize}
    \item la \textbf{variabile indipendente} è il singolo \textbf{defect};
    \item la \textbf{variabile dipendente} è la \textbf{percentuale di AV disponibili} e di AV \textbf{disponibili e consistenti}.
\end{itemize}

\subsection{Risultati}

I risultati mostrano che il problema è molto rilevante. Sul totale dei difetti selezionati:

\begin{itemize}
    \item il numero complessivo di bug chiusi e collegati a Git nei 212 progetti Apache è pari a \textbf{125.860};
    \item \textbf{63.539 bug}, cioè circa il \textbf{51\%}, risultano privi di AV utilizzabile oppure presentano un'AV inaffidabile.
\end{itemize}

Inoltre, per la maggior parte dei progetti, più del \textbf{25\%} dei difetti non possiede una AV disponibile e consistente.

\subsection{Discussione e implicazioni pratiche}

La conclusione di RQ1 è netta: l'approccio basato esclusivamente sulla AV inserita in Jira \textbf{non è sufficiente} nella maggioranza dei casi.

Se più della metà dei bug non ha una AV affidabile, allora il labeling realistico basato solo sui dati espliciti del ticket diventa spesso impossibile. Questo giustifica la necessità di metodi automatici alternativi.

\begin{notaprof}
Se non abbiamo un'AV affidabile e non introduciamo un metodo compensativo, finiamo per etichettare come pulite molte classi che invece erano già affette da difetti. Questo produce inevitabilmente molti \textbf{falsi negativi} nel dataset.
\end{notaprof}

\section{RQ2: Do methods have different accuracy in AV labeling?}

\subsection{Razionale}

Dato che RQ1 mostra che l'AV è spesso assente o inconsistente, RQ2 si chiede se esistano metodi automatici più accurati per \textbf{predire correttamente le versioni affette da un bug}.

La domanda non è più semplicemente ``l'AV c'è?'', ma: \emph{quando manca, quale metodo ricostruisce meglio l'intervallo delle versioni difettose?}

\subsection{Selezione dei progetti}

Per valutare i metodi in modo affidabile, RQ2 restringe l'analisi a progetti per cui esiste una ground truth più solida. I criteri di selezione sono:

\begin{itemize}
    \item progetti con più del \textbf{50\%} di AV disponibili e consistenti;
    \item progetti con almeno \textbf{100 bug} collegati a un Git fix commit e con AV disponibile e consistente;
    \item progetti con almeno \textbf{6 versioni}.
\end{itemize}

Dopo questa selezione, rimangono \textbf{76 progetti Apache Jira}.

\subsection{Formula della Proportion e significato di \(P\)}

Il paper introduce una misura proporzionale del ciclo di vita del difetto:

\[
P = \frac{FV - IV}{FV - OV}
\]

Questa quantità misura, in termini di versioni, il rapporto tra:

\begin{itemize}
    \item la distanza tra \textbf{Fixed Version} e \textbf{Injected Version};
    \item la distanza tra \textbf{Fixed Version} e \textbf{Opening Version}.
\end{itemize}

Una volta stimata \(P\), è possibile ricavare una \textbf{Predicted IV} mediante la formula:

\[
\text{Predicted IV} = FV - (FV - OV)\cdot P
\]

\begin{notaprof}
L'idea intuitiva è che, pur cambiando la durata assoluta dei singoli bug, la \textbf{proporzione} tra momento di scoperta e momento di correzione tenda a essere relativamente stabile. Per questo motivo si prova a usare \(P\) come ``firma temporale'' del lifecycle dei difetti.
\end{notaprof}

\subsection{Varianti di Proportion}

Il paper confronta tre modi diversi di stimare \(P\):

\begin{itemize}
    \item \textbf{Cold Start}: \(P\) è calcolato come mediana dei difetti osservati negli altri progetti;
    \item \textbf{Increment}: \(P\) è calcolato come media dei difetti corretti nelle versioni precedenti dello stesso progetto;
    \item \textbf{Moving Window}: \(P\) è calcolato come media dell'ultimo \(1\%\) dei difetti corretti.
\end{itemize}

Le tre varianti riflettono tre strategie diverse: usare esperienza di altri progetti, usare solo il passato interno del progetto, oppure concentrarsi sul comportamento più recente del processo.

\begin{notaprof}
A lezione è stata discussa anche una variante chiamata \textbf{Total}, che usa tutti i difetti passati e futuri per stimare \(P\). Questa variante non appartiene al paper, ma è utile per capire il problema del realismo temporale.
\end{notaprof}

\begin{notaesame}
La variante \textbf{Total} è concettualmente scorretta in ottica di Machine Learning realistico, perché usa informazioni del futuro per etichettare il passato. In un vero scenario predittivo, il futuro non è disponibile al momento della costruzione del dataset.
\end{notaesame}

\subsection{Metodi confrontati}

I metodi messi a confronto in RQ2 sono:

\begin{itemize}
    \item \textbf{Simple}: assume semplicemente \(IV = OV\);
    \item \textbf{SZZ\_B} (\emph{Basic});
    \item \textbf{SZZ\_U} (\emph{Unleashed});
    \item \textbf{SZZ\_RA} (\emph{Refactoring Aware});
    \item \textbf{Proportion\_ColdStart};
    \item \textbf{Proportion\_Increment};
    \item \textbf{Proportion\_MovingWindow};
    \item le versioni \textbf{SZZ\_X+}, cioè ciascun metodo SZZ combinato con \textbf{Simple}.
\end{itemize}

\subsection{Variabili dipendenti e definizione di TP, TN, FP, FN}

L'unità di analisi di RQ2 non è ancora la classe-versione, ma la \textbf{versione rispetto a un singolo difetto}. In altre parole, si valuta se un metodo etichetti correttamente una certa versione come affetta